\subsection{累積達成と到達段階の内訳}
表\ref{tab:levels}はL0〜L3の累積達成件数，表\ref{tab:exclusive}は最終到達段階の排他的内訳である．実験1は主分析，実験2は追加分析を表す．各達成率の分母は同じ行の全件数とし，L1通過例のみを分母とする条件付き正答率とは区別する．排他的な5区分の合計は全件数に一致する．モデル・条件別の全表はREADMEから参照できる固定CSVに保存した．本データではL2とL3の件数が一致し，L2のみ達成した例は0件だった．これは定義上の同一性や構造検証の十分性を示さない．
\begin{table}[t]
\centering
\caption{水準別の累積達成件数}
\label{tab:levels}
\begin{tabular}{ll|rrrrr}
\hline
実験 & 条件 & 全件 & L0 & L1 & L2 & L3 \\
\hline
1 & B & 300 & 218 & 170 & 139 & 139 \\
1 & R & 300 & 239 & 207 & 188 & 188 \\
1 & F & 300 & 256 & 244 & 219 & 219 \\
1 & D & 300 & 262 & 244 & 225 & 225 \\
1 & P & 300 & 213 & 157 & 68 & 68 \\
2 & B & 500 & 322 & 268 & 212 & 212 \\
2 & R & 500 & 343 & 303 & 265 & 265 \\
2 & F & 500 & 371 & 350 & 309 & 309 \\
2 & D & 500 & 373 & 349 & 305 & 305 \\
2 & P & 500 & 339 & 278 & 151 & 151 \\
\hline
\end{tabular}
\end{table}
\begin{table}[t]
\centering
\caption{排他的な最終到達段階の件数}
\label{tab:exclusive}
\setlength{\tabcolsep}{5pt}
\begin{tabular}{ll|rrrrr}
\hline
実験 & 条件 & 未達 & L0のみ & L1まで & L2まで & L3 \\
\hline
1 & B & 82 & 48 & 31 & 0 & 139 \\
1 & R & 61 & 32 & 19 & 0 & 188 \\
1 & F & 44 & 12 & 25 & 0 & 219 \\
1 & D & 38 & 18 & 19 & 0 & 225 \\
1 & P & 87 & 56 & 89 & 0 & 68 \\
2 & B & 178 & 54 & 56 & 0 & 212 \\
2 & R & 157 & 40 & 38 & 0 & 265 \\
2 & F & 129 & 21 & 41 & 0 & 309 \\
2 & D & 127 & 24 & 44 & 0 & 305 \\
2 & P & 161 & 61 & 127 & 0 & 151 \\
\hline
\end{tabular}
\end{table}
\subsection{対応付きの段階移動と初期成功}
表\ref{tab:movement}はbaselineから修正条件の最終評価への到達段階の移動を示す．上昇・維持・低下は順序カテゴリの比較であり，段階差を等間隔の点数として平均しない．部分改善は，上昇したもののL3未達の例で，上昇件数の内数である．主分析ではR・F・Dの部分改善が11・20・22件，追加分析では16・20・25件あった．L3救出だけではこの改善は見えない．baseline段階別の5×5遷移表も固定CSVに保存した．

主分析のbaseline L3成功139件と追加分析の212件は，R・F・Dの各条件で修正せず終了した．修正対象はそれぞれ161件と288件である．L3救出は主分析でR 49/161，F 80/161，D 86/161，追加分析でR 53/288，F 97/288，D 93/288となった．段階低下はL3成功からの退行に限らず，L1からL0未達への低下などを含む．本データでは修正型の最終L3退行は0件だが，段階低下は存在した．初期成功を含む全件達成率と，初期失敗を分母とする救出率は別の指標である．

Pは別プロンプトの生成前介入であり，修正による移動とは分けて扱う．主分析の対応付き比較では上昇47・維持141・低下112件，追加分析では96・266・138件である．L3救出は22/161件と46/288件，baseline L3成功からの退行は93/139件と107/212件だった．Pの改善例を同一コードの修正成功とは解釈しない．
\begin{table}[t]
\centering
\caption{修正条件の最終到達段階の移動}
\label{tab:movement}
\begin{tabular}{ll|rrrr}
\hline
実験 & 条件 & 上昇 & 維持 & 低下 & 部分改善 \\
\hline
1 & R & 60 & 229 & 11 & 11 \\
1 & F & 100 & 185 & 15 & 20 \\
1 & D & 108 & 177 & 15 & 22 \\
2 & R & 69 & 416 & 15 & 16 \\
2 & F & 117 & 374 & 9 & 20 \\
2 & D & 118 & 371 & 11 & 25 \\
\hline
\end{tabular}
\end{table}

\subsection{構造検証の範囲}
表\ref{tab:structure}に保存時のタスク定義に基づく構造要件を示す．構造評価は量子ビット数と測定操作を除いた回路命令数の範囲を確認する．条件分岐などの意味や個々のゲート種・接続・深さ・最適性を完全に検証するものではない．本実験の10タスクでは古典ビット数・測定数・回路内測定の必須／禁止は未指定であり，独立した構造制約としては要求しない．分布評価時に評価用コピーへ測定を追加することと，生成回路の構造制約を区別する．L3は利用可能な要件の範囲での総合達成を表す．タスク定義が実行時コミットと一致することを確認し，ファイルhashとしきい値を固定出典に保存した．
\begin{table}[t]
\centering
\caption{タスク別の構造要件}
\label{tab:structure}
\begin{tabular}{l|rr}
\hline
タスク & 量子ビット数 & 回路命令数の範囲 \\
\hline
T1 & 2 & 2--5 \\
T2 & 3 & 3--6 \\
T3a & 3 & 5--10 \\
T3b & 3 & 7--14 \\
T4 & 4 & 7--16 \\
T5 & 2 & 8--18 \\
T6 & 3 & 1--50 \\
T7 & 3 & 1--60 \\
T8 & 4 & 8--80 \\
T9 & 2 & 4--10 \\
\hline
\end{tabular}
\end{table}
